\documentclass[11pt,a4paper]{article}

\usepackage[utf8]{inputenc}
\usepackage[T1]{fontenc}
\usepackage{lmodern}
\usepackage{microtype}
% keep "--" in options as two hyphens
\DisableLigatures[-]{family=tt*}
\usepackage[margin=2.5cm]{geometry}
\usepackage{booktabs}
\usepackage{tabularx}
\usepackage{xcolor}
\usepackage{listings}
\usepackage[hidelinks]{hyperref}

\lstset{
  basicstyle=\ttfamily\small,
  backgroundcolor=\color{black!4},
  frame=single,
  rulecolor=\color{black!20},
  columns=fullflexible,
  keepspaces=true,
  breaklines=true,
  showstringspaces=false,
  xleftmargin=0.5em,
  xrightmargin=0.5em,
}

\newcommand{\cm}{\textsc{CheckMATE}}
\newcommand{\combine}{\textsc{Combine}}
\newcommand{\file}[1]{\texttt{#1}}
\emergencystretch=2em
\newcommand{\opt}[1]{\texttt{#1}}

\title{\cm{} implementation of CMS-NPS-025-002 (\file{cms\_2605\_13614})\\[0.3em]
\large 13 and 13.6\,TeV analyses and the statistical combination with \combine{}}
\author{Krzysztof Rolbiecki}
\date{October 2026}

\begin{document}
\maketitle
\tableofcontents

%======================================================================
\section{Overview}
%======================================================================

The CMS search for pair production of additional neutral scalars of the
Inert Doublet Model (IDM) in final states with two electrons or two muons
and missing transverse momentum uses the full Run~2 (2016--2018,
$\sqrt{s}=13$\,TeV) and the 2022 (Run~3, $\sqrt{s}=13.6$\,TeV) data sets.
The signal is extracted with a parametrised neural network (pNN) whose
output is fitted in the full CMS statistical model with \combine{}.

In \cm{} the analysis is split into two implementations, one per
centre-of-mass energy, and the results of the two \cm{} runs are combined
afterwards:

\begin{center}
\begin{tabular}{llll}
\toprule
Analysis & $\sqrt{s}$ & Data sets & Detector setup (experiment) \\
\midrule
\file{cms\_2605\_13614}      & 13\,TeV   & 2016, 2017, 2018 & \file{cms13tev} \\
\file{cms\_2605\_13614\_13p6} & 13.6\,TeV & 2022             & \file{cms13\_6tev} \\
\bottomrule
\end{tabular}
\end{center}

The full workflow is:
\begin{enumerate}
  \item generate signal events at 13 and 13.6\,TeV,
  \item run \cm{} on the 13\,TeV events with the analysis
        \file{cms\_2605\_13614} and on the 13.6\,TeV events with the
        analysis \file{cms\_2605\_13614\_13p6},
  \item run \file{bin/combine\_cms\_2605\_13614.py} on the two runs to
        obtain the limits for each IDM mass point.
\end{enumerate}

%======================================================================
\section{Requirements}
%======================================================================

\begin{itemize}
  \item \cm{} compiled with ONNX Runtime support (\texttt{HAVE\_ONNX}); the
        pNN model is read from
        \file{data/cms\_2605\_13614/pnn\_model\_with\_preprocessing.onnx}.
  \item The CMS datacards in \file{data/cms\_2605\_13614/datacards}:
        \begin{itemize}
          \item the main datacard \file{input/combined\_cards.txt},
          \item the histograms of every mass point,\\
                \file{input/histograms/mH<mH>\_mA<mA>/histograms.root}\\
                (about 17\,MB per mass point, 9\,GB in total).
        \end{itemize}
  \item PyROOT, used to update the signal histograms.
  \item \combine{} available in the \texttt{PATH} (the CMS results were
        produced with version~v10.5.0). It may be necessary to run
        \texttt{ulimit -s unlimited; ulimit -u unlimited} before running
        \combine{}.
\end{itemize}

%======================================================================
\section{The CMS 13.6\,TeV detector setup}
%======================================================================

The analysis handler \file{AnalysisHandlerCMS\_13\_6TeV}
(\file{tools/fritz/src/analysishandler}) provides the detector setup for
CMS analyses at $\sqrt{s}=13.6$\,TeV. All its settings, i.e.\ the photon,
electron and muon identification efficiencies and the $b$- and $\tau$-tagging
efficiencies, are taken over unchanged from the 13\,TeV handler
\file{AnalysisHandlerCMS\_13TeV}. The Delphes card of the 13\,TeV setup,
\file{delphes\_skimmed\_CMS\_13TeV.tcl}, is used as well.

The setup is registered in \cm{} as
\begin{center}
\begin{tabular}{ll}
\toprule
Experiment                 & \file{cms13\_6tev} \\
Analysis group             & \file{CMS\_13\_6TeV} \\
Centre-of-mass energy      & 13.6\,TeV \\
Analysis list              & \file{data/analysis\_info/CMS\_13\_6TeV\_analyses.txt} \\
Analysis sources           & \file{tools/analysis/\{src,include\}/analyses/CMS\_13\_6TeV} \\
\bottomrule
\end{tabular}
\end{center}

New 13.6\,TeV analyses are added with the \file{AnalysisManager} by
choosing CMS and a centre-of-mass energy of 13.6\,TeV. The experiment
recorded in the analysis settings file (\file{\_var.j}) selects the
analysis handler at run time:
\begin{lstlisting}
"ecm": "13.6",
"experiment": "cms13_6tev",
\end{lstlisting}

\paragraph{Note.} Analysis names must be unique in \cm{} across all
analysis groups. A 13.6\,TeV version of an existing analysis therefore needs
a new name, e.g.\ with the suffix \file{\_13p6}. When the settings file of
an existing analysis is copied by hand, the \texttt{experiment} entry has to
be changed as well; otherwise the analysis is booked in the 13\,TeV handler
and \cm{} stops with \emph{``Cannot load analysis \ldots{} -- analysis
unknown''}.

%======================================================================
\section{The analyses}
%======================================================================

\subsection{Event selection and pNN evaluation}

Both implementations apply the event preselection of the CMS analysis to
events with two same-flavour leptons and evaluate the pNN for every mass
point $(m_H, m_A)$ of the mass grid (Section~\ref{sec:grid}). Events with a
pNN output of at least 0.9 are counted in the bins of the pNN distribution
used by CMS: 7 bins in the electron channel and 11 bins in the muon channel.

\begin{itemize}
  \item \file{cms\_2605\_13614} (13\,TeV, $\mathcal{L}=138$\,fb$^{-1}$):
        every event is assigned at random to one of the data-taking years
        2016, 2017 and 2018, with probabilities 25.3\,\%, 27.8\,\% and
        46.9\,\% according to the luminosities of the years. The
        year-dependent lepton trigger thresholds are applied.
  \item \file{cms\_2605\_13614\_13p6} (13.6\,TeV, $\mathcal{L}=35$\,fb$^{-1}$):
        all events are assigned to 2022.
\end{itemize}

\subsection{Output: control regions}

The yields are stored as control regions in the \cm{} output file
\begin{center}
\file{results/<run>/analysis/<process>\_<analysis>\_control.dat}
\end{center}
with one control region per mass point, year, channel and pNN bin:
\begin{center}
\file{mH<mH>\_mA<mA>\_<year>\_<channel>\_bin\_<n>}
\end{center}
where \file{<channel>} is \file{EE} (bins \file{0}--\file{6}) or \file{MM}
(bins \file{00}--\file{10}), e.g.\ \file{mH70.0\_mA140.0\_2016\_EE\_bin\_3}.
Both analyses book the regions of all four years, so that the output files
of the two energies have the same structure; the regions of the years not
covered by an analysis remain empty.

\subsection{Mass grid}
\label{sec:grid}

The mass points are defined in the \texttt{mHmA} list in the analysis
headers \file{cms\_2605\_13614.h} and \file{cms\_2605\_13614\_13p6.h}. The
list contains all 524 mass points for which CMS provides histograms in the
datacards:
\begin{itemize}
  \item $m_H = 60$--$150$\,GeV in steps of 2.5\,GeV,
  \item $m_A$ mostly on a 5\,GeV grid with
        $20\,\mathrm{GeV} \le m_A - m_H \le 100\,\mathrm{GeV}$, including
        additional points at $m_A = m_H+92$\,GeV; for $m_H \ge 135$\,GeV the
        grid is $m_A = m_H+20, 30, \ldots, 100$\,GeV and $m_H+92$\,GeV.
\end{itemize}
The exact list of points is given in the headers.
Since the control regions are cheap to produce, \cm{} always evaluates the
full grid. The selection of the mass points for the statistical evaluation
is done in the combination script.

%======================================================================
\section{Running \cm{}}
%======================================================================

The two energies are run separately, each with events generated at the
corresponding centre-of-mass energy. Example parameter files:

\begin{lstlisting}
[Parameters]
Name: wino_320_40_decay
Analyses: cms_2605_13614

[hhh]
Events: /path/to/events_13TeV.hepmc
XSect: 18.8 fb
\end{lstlisting}

\begin{lstlisting}
[Parameters]
Name: wino_320_40_decay_136
Analyses: cms_2605_13614_13p6

[hhh]
Events: /path/to/events_13p6TeV.hepmc
XSect: 20.8 fb
\end{lstlisting}

All control-region files of a run (one per process) are taken into account
in the combination.

%======================================================================
\section{Statistical combination: \file{combine\_cms\_2605\_13614.py}}
%======================================================================

\subsection{Synopsis}

\begin{lstlisting}
bin/combine_cms_2605_13614.py run [run ...]
    [--mH MH --mA MA | --all]
    [--datacards DIR] [--output DIR]
    [--r-max R] [--combine-executable EXE]
    [--prepare-only] [--keep-histograms]
\end{lstlisting}

The script is generated by \texttt{configure} from
\file{tools/python/combine\_cms\_2605\_13614.in}, in the same way as
\file{bin/CheckMATE}. It therefore uses the same results and data folders as
\cm{}: the results folder is \file{results/} in the \cm{} directory, or the
folder given in the environment variable \texttt{CHECKMATE\_RESULTS}.

\subsection{Description}

The script
\begin{enumerate}
  \item collects the files \file{*\_cms\_2605\_13614*\_control.dat} from
        \file{<results>/<run>/analysis/} of all given runs, i.e.\ the output
        of both \file{cms\_2605\_13614} and \file{cms\_2605\_13614\_13p6};
  \item sums the yields of every control region over all files. The
        statistical uncertainty of each file,
        $\Delta N = N / \sqrt{A \cdot N_\mathrm{MC}}$ with the acceptance $A$
        and the number of generated events $N_\mathrm{MC}$, is added in
        quadrature. The result is written to
        \file{cms\_2605\_13614\_combined\_control.txt} in the output folder;
  \item evaluates the selected mass points one after another. For every
        mass point it
        \begin{enumerate}
          \item copies the original \file{histograms.root} of the mass point
                from the datacards to the output folder,
          \item replaces the signal histograms of the copy (the
                \texttt{nominal} histogram and all systematic variations in
                \file{<year>/SR/\{Electron,Muon\}/idm}) by the combined yields,
          \item runs
\begin{lstlisting}
combine combined_cards.txt -M AsymptoticLimits \
    --keyword-value MA=<mA> --keyword-value MH=<mH> \
    --strictBounds --rMax <R> -n _mH<mH>_mA<mA>
\end{lstlisting}
                in the output folder, which contains a copy of
                \file{combined\_cards.txt},
          \item extracts the observed and expected limits on the signal
                strength $r$ and adds them to \file{limits.txt},
          \item deletes the modified histogram file.
        \end{enumerate}
\end{enumerate}

\paragraph{Handling of the histograms.} The datacards are never modified.
Only one modified histogram file (about 17\,MB) exists at any time, and it
is deleted after the mass point has been evaluated, also when the evaluation
fails or the script is interrupted (Ctrl-C). Modified histograms are only
kept with \opt{--keep-histograms} or \opt{--prepare-only}; note that this
requires about 1.5\,GB for the default selection and 9\,GB with
\opt{--all}.

\subsection{Arguments and options}

\begin{center}
\begin{tabularx}{\textwidth}{lX}
\toprule
\opt{run [run ...]} & Names of the \cm{} runs to combine (folders in the
                      results folder), typically the 13\,TeV and the
                      13.6\,TeV run. A single run can be given as well. \\
\opt{--mH MH --mA MA} & Evaluate a single mass point. Both masses have to be
                      given, and the point has to be in the mass grid of
                      the analysis and have histograms in the datacards. \\
\opt{--all}         & Evaluate all 524 mass points of the analysis. \\
\opt{--datacards DIR} & Folder with the \combine{} datacards; default:\newline
                      \file{data/cms\_2605\_13614/datacards}. \\
\opt{--output DIR}  & Output folder (default:
                      \file{combined\_<run1>\_<run2>\ldots} in the results
                      folder). \\
\opt{--r-max R}     & Maximum signal strength passed to \combine{}
                      (default: 10). \\
\opt{--combine-executable EXE} & \combine{} executable (default:
                      \texttt{combine}). \\
\opt{--prepare-only} & Only prepare the histograms, do not run \combine{}.
                      The modified histograms are kept. \\
\opt{--keep-histograms} & Keep the modified histograms after running
                      \combine{}. \\
\bottomrule
\end{tabularx}
\end{center}

\subsection{Selection of mass points}

\begin{center}
\begin{tabular}{ll}
\toprule
Command line & Evaluated mass points \\
\midrule
(default)              & 90 points with $m_H$ and $m_A$ on a 10\,GeV grid \\
\opt{--all}            & all 524 points \\
\opt{--mH X --mA Y}    & the single point $(X, Y)$, any of the 524 points \\
\bottomrule
\end{tabular}
\end{center}

\opt{--all} cannot be combined with \opt{--mH}/\opt{--mA}. Mass points
without histograms in the datacards are skipped with a warning.

\subsection{Output}

The output folder contains
\begin{description}
  \item[\file{cms\_2605\_13614\_combined\_control.txt}] combined
        control-region yields and uncertainties of all mass points,
  \item[\file{combined\_cards.txt}] copy of the \combine{} datacard,
  \item[\file{limits.txt}] table of the limits on $r$ per mass point (see
        below),
  \item[\file{logs/mH<mH>\_mA<mA>.log}] full \combine{} output of every
        mass point,
  \item[\file{higgsCombine\_mH<mH>\_mA<mA>.AsymptoticLimits.mH120.root}]
        \mbox{}\\ \combine{} output file of every mass point.
\end{description}

\file{limits.txt} has the columns
\begin{lstlisting}
mH  mA  obs  exp2.5  exp16  exp50  exp84  exp97.5  status
\end{lstlisting}
with the observed limit, the expected limit and its quantiles, and the
status \texttt{ok}, \texttt{failed} or \texttt{prepared}
(\opt{--prepare-only}); missing values are written as \texttt{-}. A mass
point is excluded at 95\,\%~CL if $r_\mathrm{obs} < 1$. Results of earlier
calls with the same output folder are kept, so the table can be filled
point by point. The script ends with exit code 1 if any mass point failed
and lists these points.

\subsection{Examples}

Evaluate a single mass point:
\begin{lstlisting}
bin/combine_cms_2605_13614.py wino_320_40_decay wino_320_40_decay_136 \
    --mH 70 --mA 140
\end{lstlisting}

Evaluate the default selection of 90 mass points:
\begin{lstlisting}
bin/combine_cms_2605_13614.py wino_320_40_decay wino_320_40_decay_136
\end{lstlisting}

Evaluate all mass points:
\begin{lstlisting}
bin/combine_cms_2605_13614.py wino_320_40_decay wino_320_40_decay_136 --all
\end{lstlisting}

Prepare the histograms of one mass point for inspection without running
\combine{}:
\begin{lstlisting}
bin/combine_cms_2605_13614.py wino_320_40_decay wino_320_40_decay_136 \
    --mH 70 --mA 140 --prepare-only --output /tmp/check
\end{lstlisting}

\subsection{Helper scripts}

The script uses two modules in \file{tools/python} that can also be run on
their own:
\begin{description}
  \item[\file{combine\_control.py <folder> <analysis>}] \mbox{}\\
        sums the control-region files \file{*\_<analysis>\_control.dat}
        in one folder and writes the result to
        \file{<analysis>\_combined\_control.txt}.
  \item[\file{combine\_wrapper.py <mH> <mA> --control-file F --histogram-file H}]
        \mbox{}\\ fills the combined yields of one mass point into a
        histogram file. \emph{This modifies the given file in place.}
\end{description}

\end{document}
